\section{Memory Management Framework}
\label{sec:memory}

\subsection{The 28GB Constraint}

The GMKtech M7's 28GB total memory (with 4GB allocated to iGPU VRAM) represents a challenging constraint for LLM deployment. This section details our memory management framework designed to operate effectively within this limit.

\subsubsection{Memory Budget Analysis}

Breaking down the available 24GB for model inference:

\begin{table}[h]
\centering
\caption{Memory Budget Allocation}
\label{tab:memory_budget}
\begin{tabular}{llc}
\toprule
\textbf{Component} & \textbf{Description} & \textbf{Allocation (GB)} \\
\midrule
OS / WSL2 & System overhead & 4.0 (already reserved) \\
Working Memory & Temporary buffers & 2.0 \\
Safety Margin & Headroom for spikes & 2.0 \\
Model Weights & Quantized parameters & 1.5 - 2.5 \\
KV Cache & Attention state storage & 14.0 - 16.0 \\
\midrule
\textbf{Total} & & \textbf{19.5 - 24.5} \\
\bottomrule
\end{tabular}
\end{table}

The KV cache dominates memory consumption, motivating our aggressive compression strategy.

\subsection{Automatic Quantization Selection}

The MemoryManager automatically selects appropriate quantization based on available memory and model size.

\subsubsection{Selection Algorithm}

The quantization selection process considers:

\begin{enumerate}
    \item \textbf{Model Size:} Parameters in billions
    \item \textbf{Target Context:} Desired maximum context length
    \item \textbf{Available Memory:} Current system state
    \item \textbf{Quality Requirements:} Minimum acceptable accuracy
\end{enumerate}

\begin{algorithm}[H]
\caption{Quantization Selection Algorithm}
\label{alg:quant_selection}
\begin{algorithmic}[1]
\Require Model parameters $P$ (billions), Available memory $M$ (GB)
\Ensure Selected quantization $Q$
\State $S_{FP16} \gets P \times 2$ \Comment{FP16 size in GB}
\State $ratios \gets \{Q8\_0: 1.0, Q6\_K: 0.75, Q5\_K\_M: 0.625, Q4\_K\_M: 0.5, Q3\_K\_M: 0.375\}$
\State $KV_{overhead} \gets 2.6$ \Comment{Approximate for 192K context with turbo3}
\State $working \gets 2.0$
\For{$Q$ in preference\_order}
    \State $model\_size \gets S_{FP16} \times ratios[Q]$
    \State $total \gets model\_size + KV_{overhead} + working$
    \If{$total \leq M \times 0.85$}
        \State \Return $Q$
    \EndIf
\EndFor
\State \Return $Q2\_K$ \Comment{Fallback to most aggressive}
\end{algorithmic}
\end{algorithm}

\subsubsection{Preference Ordering}

The default preference ordering (Q4\_K\_M, Q5\_K\_M, Q6\_K, Q8\_0, Q3\_K\_M, Q2\_K) reflects a balance between quality and compression:

\begin{itemize}
    \item Q4\_K\_M offers the best quality/size trade-off for most applications
    \item Q5\_K\_M for quality-critical applications with memory headroom
    \item Q6\_K and Q8\_0 when quality is paramount
    \item Q3\_K\_M and Q2\_K only when memory is severely constrained
\end{itemize}

\subsection{Dynamic Context Window Calculation}

Beyond model size, the context window significantly impacts memory requirements.

\subsubsection{KV Cache Size Formula}

For a given context length $T$ and KV cache quantization $Q$:

\begin{equation}
M_{KV}(T, Q) = 2 \times L \times H \times d_h \times T \times b(Q)
\end{equation}

Where:
\begin{itemize}
    \item $L = 32$ (number of layers)
    \item $H = 32$ (number of heads)
    \item $d_h = 96$ (head dimension)
    \item $b(Q)$ is bytes per element for quantization $Q$
\end{itemize}

Bytes per element for different quantization schemes:

\begin{table}[h]
\centering
\caption{KV Cache Bytes per Element by Quantization}
\label{tab:kv_bytes}
\begin{tabular}{lc}
\toprule
\textbf{Quantization} & \textbf{Bytes per Element} \\
\midrule
FP16 & 2.0 \\
Q8\_0 & 1.0 \\
Q4\_0 & 0.5 \\
Turbo4 & 0.5625 ($\approx$ 4.5 bits) \\
Turbo3 & 0.4375 ($\approx$ 3.5 bits) \\
\bottomrule
\end{tabular}
\end{table}

\subsubsection{Context Size Calculation}

Given available memory $M_{avail}$ after model and working memory:

\begin{equation}
T_{max} = \left\lfloor \frac{M_{avail} \times 1024^3}{2 \times L \times H \times d_h \times b(Q)} \right\rfloor
\end{equation}

For MiniMax M2.7 with turbo3 (14KB/token effective):

\begin{itemize}
    \item With 6GB available: $T_{max} \approx 458,000$ tokens (theoretical)
    \item Practical limit: 194,560 (192K usable after safety margin)
\end{itemize}

\subsection{Safety Margins and Limits}

The memory manager implements multiple safety mechanisms:

\subsubsection{Safety Factor}

Default safety factor of 0.85 ensures headroom:

\begin{equation}
M_{usable} = M_{available} \times 0.85
\end{equation}

This accommodates:
\begin{itemize}
    \item Temporary allocations during generation
    \item Memory fragmentation
    \item Operating system fluctuations
    \item Garbage collection pauses (Python)
\end{itemize}

\subsubsection{Hard Limits}

Hard limits prevent allocation beyond physical capacity:

\begin{lstlisting}[language=Python, caption=Memory Limit Enforcement]
def check_memory_available(self, required_gb: float) -> bool:
    stats = self.get_stats()
    # Account for already allocated model/cache memory
    available = (
        stats.available_gb - 
        self.current_model_memory - 
        self.current_kv_memory
    )
    return available >= required_gb
\end{lstlisting}

\subsection{Memory Tracking and Statistics}

The MemoryManager provides comprehensive memory tracking:

\begin{lstlisting}[language=Python, caption=Memory Statistics]
@dataclass
class MemoryStats:
    total_gb: float          # Total system RAM
    available_gb: float      # Currently available
    used_gb: float          # Currently used
    percent_used: float     # Usage percentage
    model_memory_gb: float  # Model weights
    kv_cache_gb: float      # KV cache allocation
\end{lstlisting}

\subsection{Optimal Settings Generation}

The \texttt{get\_optimal\_settings} method generates complete configurations:

\begin{lstlisting}[language=Python, caption=Optimal Settings Generation]
def get_optimal_settings(
    self,
    model_params_billions: float
) -> Dict[str, Any]:
    quant = self.select_quantization(model_params_billions)
    
    # Calculate model size
    fp16_size = model_params_billions * 2
    ratios = {"Q8_0": 1.0, "Q6_K": 0.75, "Q5_K_M": 0.625,
              "Q4_K_M": 0.5, "Q3_K_M": 0.375, "Q2_K": 0.25}
    model_size = fp16_size * ratios.get(quant, 0.5)
    
    # Context and batch sizing
    kv_type = "turbo3"
    max_ctx = self.calculate_max_context(model_size, kv_type)
    
    cpu_count = psutil.cpu_count(logical=False) or 4
    
    # Batch sizes based on context
    if max_ctx >= 131072:
        batch_size = 2048
        ubatch_size = 512
    elif max_ctx >= 8192:
        batch_size = 1024
        ubatch_size = 512
    else:
        batch_size = 512
        ubatch_size = 256
    
    return {
        "quantization": quant,
        "n_ctx": max_ctx,
        "cache_type_k": kv_type,
        "cache_type_v": kv_type,
        "n_threads": cpu_count,
        "n_batch": batch_size,
        "n_ubatch": ubatch_size,
        "flash_attn": True,
    }
\end{lstlisting}

\subsection{Context Fallback Strategy}

When initial allocation fails, the system implements context fallback:

\begin{enumerate}
    \item Attempt requested context size
    \item If OOM, try progressively smaller sizes: 131K, 98K, 65K, 49K, 32K, 16K, 8K, 4K
    \item Report actual context achieved
    \item Allow user to accept or retry with different quantization
\end{enumerate}

This graceful degradation ensures the model loads successfully even under memory pressure.

\subsection{Working Memory Reservation}

Working memory accommodates temporary allocations during generation:

\begin{itemize}
    \item \textbf{Activation buffers:} Intermediate computation results
    \item \textbf{Attention matrices:} Temporary attention scores
    \item \textbf{Sampling state:} Temperature, top-k/p tracking
    \item \textbf{Python overhead:} GC, interpreter overhead
\end{itemize}

We reserve 2GB for working memory, sufficient for MiniMax M2.7 with substantial headroom.

\subsection{Memory Profiling and Diagnostics}

The framework includes diagnostic capabilities:

\begin{lstlisting}[language=Python, caption=Memory Profiling]
async def profile_memory_usage(self, prompt: str, max_tokens: int):
    tracemalloc.start()
    
    # Baseline
    baseline = self.get_memory_info()
    
    # Generate
    result = await self.model.generate(prompt, max_tokens=max_tokens)
    
    # Peak
    current, peak = tracemalloc.get_traced_memory()
    tracemalloc.stop()
    
    return {
        "baseline_mb": baseline.rss_mb,
        "peak_mb": peak / 1024 / 1024,
        "increase_mb": (peak - baseline.rss * 1024 * 1024) / 1024 / 1024,
        "snapshots": tracemalloc.snapshot(),
    }
\end{lstlisting}

\subsection{Memory Efficiency Score}

We define a composite memory efficiency score:

\begin{equation}
\text{Efficiency} = 0.4 \times \frac{M_{theoretical}}{M_{actual}} + 0.3 \times \left(1 - \frac{M_{peak} - M_{post\_gc}}{M_{peak}}\right) + 0.3 \times \frac{T_{achieved}}{T_{target}}
\end{equation}

Where:
\begin{itemize}
    \item Theoretical vs actual memory: How close to optimal allocation
    \item GC recovery: How much memory is released after cleanup
    \item Context achievement: Fraction of target context achieved
\end{itemize}

Scores range from 0 to 100, with Miniforge typically achieving 85-95 on the M7 platform.
